\chapter{Einleitung}
Das \www ist in den letzten drei Jahren zum gr"o"sten Informations- speicher und -verteiler der Welt gewachsen. Dies wird vor allem durch die zunehmende Akzeptanz in der breiten "Offentlichkeit gef"ordert. Diese Akzeptanz hat viele Firmen dazu bewegt, ihre Produkte und Informationen im \www zur Verf"ugung zu stellen.

Die Entwicklung "`alles und jedes"' im \www anzubieten, hat neben vielen seri"osen Informationsanbietern wie z.B. dem Universallexikon Encyclopaedia Britannica\footnote{http://www.eb.com/} oder dem Nachrichtendienst Reuters\footnote{http://www.reuters.com/} allerdings auch Stilbl"uten wie z.B. dem virtuellem Weiterleben\footnote{http://www.ewigesleben.de/} oder einfach nur Quatschseiten\footnote{http://www.quackwatch.com/} hervorgebracht.

Da das \www eigentlich nur f"ur die institutsinterne Verteilung von Dokumenten vorgesehen war, wurden auch keinerlei Vorkehrungen getroffen, Dokumente zu finden oder in geordneten Gruppen zusammenfassen zu k"onnen. Mit einem Datenbestand von nun mehr als 150 Gigabyte ( gesch"atzter Stand April 1997 ) wurde bei der Konzeption nicht gerechnet.

Um sich noch einigerma"sen, im immer gr"o"ser werdenden Informationsdschungel des \www zurechtzufinden wurden sehr bald sog. Suchmaschinen entwickelt. Suchmaschinen erlauben es dem Benutzer Dokumente via Schlagwortsuche zu finden.

Diese Suchmaschinen arbeiten jedoch momentan nur mit sehr rudiment"aren Strategien zur Informationsaufbereitung, weshalb ich dieser Diplomarbeit versuche, Konzepte f"ur eine Verbesserung der Suchmachschinen aufzuzeigen und zu implementieren:

\begin{itemize}
	\item	Entwurf und Implementierung eines optimierten Crawlers.
	
			Hier wird gezeigt, wie durch ein fehlertolerantes, verteiltes System die Suchgeschwindigkeit bestehender Crawler um ein vielfaches erh"oht werden kann.
	\item	Entwurf und Implementierung eines Sprachenerkenners f"ur die Indexierung.
	
			Sprachenerkennung bezieht sich hierbei auf die Sprache in der die Dokumente abgefa"st sind, die von Suchmaschinen indiziert werden. F"ur Benutzer von Suchmaschinen sind nur Dokumente interessant, die sie auch verstehen. Ist die Sprache, in der Dokumente abgefa"st sind bekannt, kann man den Benutzern ihnen verst"andliche Dokumente zeigen.
			
	\item	Entwurf einer kontextabh"angigen Indexierung.
	
			Indexierung der richtigen Inhalte liefert wesentlich bessere Suchantworten als dies mit koventioneller Wortindizierung m"oglich ist.
\end{itemize}

Um eine m"oglichst gro"se Portabilit"at der entwickelten Programme zu gew"ahrleisten, sind die in dieser Arbeit vorgestellten Algorithmen plattformunabh"angig in Perl5 und Objective-C implementiert.

Die Durchf"uhrung der Diplomarbeit fand am {\em Centrum f"ur Informations- und Sprachverarbeitung (CIS)} der Ludwig-Maximilians Universit"at M"unchen statt. 

Die Notwendigkeit der vorgestellten Verbesserungen wurde schon w"ahrend der Erstellung dieser Arbeit sichtbar: Zur Zeit werden die hier beschriebenen Methoden zur Sprachenerkennung von {\em AltaVista} in den dortigen Suchdienst integriert.


